\documentclass[11pt]{article}

\usepackage[margin=1in]{geometry}
\usepackage[T1]{fontenc}
\usepackage[utf8]{inputenc}
\usepackage{amsmath}

\title{Review of ``Objective-Free Entity Assembly in Block Worlds''}
\author{Anonymous Reviewer}
\date{March 6, 2026}

\begin{document}
\maketitle

\section*{Executive Summary}
This manuscript tackles an important objective-free ALife question with a strong experimental design: large simulation volume, dual Assembly Theory (AT) formulations, null-model auditing, parameter sweeps, and a catalytic control. A final whole-paper consistency check suggests the overall argument is coherent, but two paper-wide tensions still reduce interpretability: the no-reuse assembly metric appears to be edge-count-driven by construction, and several claims are phrased more strongly than the statistical resolution supports in a small-entity regime. The reuse-aware analysis and robustness checks remain the most informative components, and they can support a clear boundary-condition result if claims, equations, and inference units are aligned consistently across sections. With focused revisions on formal definitions, inferential framing, and tone calibration, the paper can serve as a credible negative-result baseline for follow-on objective-free work.

\section*{Prioritized Findings}
\textbf{The findings below are ranked by confidence (highest to lower confidence).}

\begin{enumerate}
\item \textbf{Very high confidence: The no-reuse equation likely makes one central result definitional.}  
Under edge-addition from isolated vertices, each edge is added once, so the recurrence appears equivalent to $a_{\text{exact}}(G)=|E(G)|$ for connected entities. If true, ``size-driven'' behavior for this metric is expected a priori rather than empirically discovered.

\item \textbf{High confidence: Null-model resolution is limited where most data live.}  
The entity distribution is dominated by very small graphs, where degree-preserving swap nulls have low state-space diversity and coarse p-value granularity. This constrains how strongly ``0\% excess'' can be interpreted.

\item \textbf{High confidence: Statistical emphasis should track the effective unit of inference.}  
Run-level and type-level checks are present but secondary in narrative emphasis. Promoting those units over per-snapshot counts would improve coherence between methods, results, and conclusions.

\item \textbf{Medium-high confidence: Reuse-aware formalism needs stricter specification.}  
The current recurrence mixes edge-set partitions with object-level language; connectivity of intermediates, typed-graph constraints, and canonicalization assumptions should be stated explicitly for auditability.

\item \textbf{Medium confidence: One cross-section identity appears inconsistent.}  
A statement tied to $|E|-1$ conflicts with calibration examples and surrounding definitions. Harmonizing this across text, equations, and tables is necessary for internal consistency.

\item \textbf{Medium confidence: Mechanism-level quantitative claims are plausible but under-supported.}  
The birth--death crossover explanation for the size ceiling is reasonable, yet estimation details, uncertainty bands, and fit diagnostics are not sufficiently reported.

\item \textbf{Medium confidence: Positive control interpretation should be sharpened.}  
The catalytic control demonstrates ecological sensitivity (more/larger entities), but it is not a structure-selective control. A motif-targeted control would better validate structural-detection sensitivity.

\item \textbf{Medium confidence: Tone and scope still need paper-wide calibration.}  
Conclusions should consistently be framed as conditional on the tested rule expressiveness, parameter ranges, and time horizon rather than generalized to objective-free systems broadly.
\end{enumerate}

\section*{Section Recommendations}

\subsection*{Narrative Coherence}
\begin{itemize}
\item Keep one through-line from introduction to conclusion: ``scoped boundary condition under low-expressiveness random rules.''
\item Distinguish clearly between mathematical consequences of definitions and empirical findings from simulation.
\item Use the same inferential unit in claim sentences and figure/table interpretations.
\end{itemize}

\subsection*{Abstract \& Introduction}
\begin{itemize}
\item Replace universal wording with conditional language anchored to the tested regime.
\item Add one concise sentence stating what is newly contributed beyond prior AT or objective-free baselines.
\item Flag early that the reuse-aware formulation carries most of the nontrivial structural signal.
\end{itemize}

\subsection*{Methods: Dynamics \& Rule Table}
\begin{itemize}
\item Define neighborhood semantics once (radius, metric, and drift/bond consistency) and reference it uniformly.
\item Specify tie-breaking, random-seed policy, and any deterministic ordering used in canonicalization.
\item Include a compact reproducibility table of constants, sampling cadence, and logging rules.
\end{itemize}

\subsection*{Methods: Assembly Formalism}
\begin{itemize}
\item State explicitly whether $a_{\text{exact}}$ reduces to $|E|$ for connected typed graphs, and cite or prove the claim.
\item Rewrite reuse-aware recursion using graph objects with connected intermediates, not only edge-set notation.
\item Add one worked micro-example that traces both formulations step-by-step.
\end{itemize}

\subsection*{Methods: Null Model \& Inference}
\begin{itemize}
\item For small $n$, supplement swap nulls with exact enumeration or explicitly bounded limitations.
\item Report p-value discreteness, tie frequency, and effective null support by size class.
\item State multiplicity handling and explain how many decisions are made at each inference layer.
\end{itemize}

\subsection*{Results \& Figures}
\begin{itemize}
\item Lead each major result with run-level or type-level estimates, then present raw observation totals.
\item Quantify ``size-driven'' with a compact effect-size summary for $a_{\text{reuse}}$.
\item Ensure every ``0\%'' or ``exact'' statement specifies denominator, test, and inference unit.
\item Separate ecological shifts from structural-specificity outcomes in captions and narrative text.
\end{itemize}

\subsection*{Discussion \& Conclusion}
\begin{itemize}
\item Frame the contribution as a baseline map of where complexity was not detected, not a universal impossibility claim.
\item Separate ``not observed here'' from ``cannot emerge,'' and keep that distinction consistent in all closing claims.
\item Prioritize falsifiable next experiments (partner-specific rules, motif-gated catalysis, adaptive rules).
\end{itemize}

\subsection*{Global Consistency}
\begin{itemize}
\item Align notation, claim strength, and inference unit in every section where assembly results are interpreted.
\item Remove any sentence that implies stronger generality than the method and null resolution can support.
\item Cross-check equations, prose, and tables so identical symbols carry identical meaning throughout.
\end{itemize}

\end{document}